\documentclass[14pt,a4paper]{extarticle}
\pagestyle{empty}
\usepackage{fontspec}
\usepackage{polyglossia}
\setmainlanguage{ukrainian}
\IfFontExistsTF{Times New Roman}{\setmainfont{Times New Roman}\newfontfamily\ukrainianfont{Times New Roman}}{\setmainfont{Libertinus Serif}\newfontfamily\ukrainianfont{Libertinus Serif}}
\IfFontExistsTF{Courier New}{\setmonofont{Courier New}\newfontfamily\ukrainianfonttt{Courier New}}{\setmonofont{DejaVu Sans Mono}\newfontfamily\ukrainianfonttt{DejaVu Sans Mono}}
\usepackage[top=20mm,bottom=20mm,left=25mm,right=15mm]{geometry}
\usepackage{setspace}
\onehalfspacing
\usepackage{indentfirst}
\setlength{\parindent}{1.25cm}
\setlength{\parskip}{0pt}
\usepackage{graphicx,float,caption,enumitem,titlesec}
\usepackage{tikz}
\usetikzlibrary{arrows.meta,positioning}
\graphicspath{{images/}{./}}
\captionsetup{font=it,justification=centering}
\setlist[enumerate]{label=\arabic*.,leftmargin=1.25cm,itemsep=0.3em}
\titleformat{\section}{\normalfont\normalsize\bfseries\centering}{\thesection.}{0.5em}{}
\titlespacing*{\section}{0pt}{1.2\baselineskip}{0.7\baselineskip}
\newcommand{\shot}[2]{%
	\begin{figure}[H]\centering
		\IfFileExists{images/#1.png}{\includegraphics[width=1\textwidth,height=1\textheight,keepaspectratio]{images/#1.png}}{%
			\IfFileExists{images/#1.jpg}{\includegraphics[width=0.4\textwidth,height=0.36\textheight,keepaspectratio]{images/#1.jpg}}{%
				\IfFileExists{images/#1.jpeg}{\includegraphics[width=0.78\textwidth,height=0.43\textheight,keepaspectratio]{images/#1.jpeg}}{%
					\fbox{\begin{minipage}[c][6cm][c]{0.85\textwidth}\centering
							Місце для скріншота\\[0.5cm]\texttt{\detokenize{images/#1.png}}\\[0.3cm]
							#2\end{minipage}}}}}
		\caption{#2}\label{fig:#1}\end{figure}}
\begin{document}
	\begin{titlepage}
		\begin{center}
			Міністерство освіти і науки України

			Національний університет «Львівська політехніка»

			\vspace{0.7cm}
			\IfFileExists{logo.png}{\includegraphics[width=0.9\textwidth,height=7cm,keepaspectratio]{logo.png}}{\fbox{\begin{minipage}[c][3cm][c]{0.4\textwidth}\centering Логотип університету\end{minipage}}}

			\vspace{1cm}
			\textbf{Звіт до лабораторної роботи №1}

			з дисципліни

			«Програмування мобільних додатків»

			на тему

			«Налаштування середовища та концепція додатку»
			\vspace{1.5cm}
		\end{center}
		\begin{flushright}
			Виконав:\\студент групи АП-41\\Чижевський М. М.\\[1cm]
			Прийняв:\\доц. каф. ІКТ, докт. ф.\\Андрущак В. С.
		\end{flushright}
		\vfill
		\begin{center}Львів 2026\end{center}
	\end{titlepage}

	\section*{Мета роботи}
Ознайомитися з Android Studio, налаштувати середовище розробки та обрати тему власного Android-додатка.

\section{Налаштування середовища}
Для роботи я використав Android Studio Rabbit 1 (2026.2.1) на Windows 11. Створив проєкт за шаблоном Empty Activity, обрав Kotlin та Jetpack Compose. Під час налаштування встановив Android SDK і компоненти для запуску емулятора.

У Device Manager налаштував віртуальний пристрій для перевірки базового проєкту. Для запуску на смартфоні увімкнув режим розробника та налагодження. Скріншоти запуску наведено нижче.
\shot{s1_lab1}{Базовий проєкт на емуляторі}
\shot{s2_lab1}{Базовий проєкт на смартфоні}

\section{Вибір галузі}
Я обрав електроніку. Під час роботи зі схемами мені часто потрібно підібрати номінал резистора, перевірити частоту зрізу фільтра або перевести одиниці. Тому вирішив зробити калькулятор, у якому ці інструменти будуть зібрані разом і працюватимуть без інтернету.

\section{Ідея додатка}
Назва додатка — \textbf{Ohmlette}. Основна версія призначена для Android, пізніше планується версія для ПК.

Користувач обирає розрахунок, вводить параметри й отримує результат із формулою. Почну з одиниць, номіналів, резисторних мереж і RC-фільтрів. Для автопідбору резисторів передбачено максимум шість компонентів із можливістю змішаного з'єднання.

Ще планую додати читання даташитів за допомогою локальної AI-моделі. Користувач завантажить PDF і вкаже артикул, а додаток складе українську картку компонента з параметрами, умовами та посиланнями на сторінки. Почну з діодів, резисторів, конденсаторів і дроселів. Модель можна буде завантажити в додатку, після чого обробка працюватиме без інтернету.

Щоб уточнити набір функцій, я обговорив ідею з AI. Фрагмент діалогу:
\begin{quote}
\textbf{Запит:} «Автопідбір резисторів: ну я думаю можна обмежити кількість компонентів до шести».

\textbf{Відповідь:} «До шести резисторів — розумна межа автопідбору. Змішані послідовно-паралельні схеми залишаємо».
\end{quote}
Я залишив функції, якими користувався б сам, а розрахунки для АЦП/ЦАП та операційних підсилювачів прибрав.

\section{Вимоги та структура додатка}
Основні функції:
\begin{enumerate}
\item Арифметичні дії та перерахунок одиниць із префіксами від піко до гіга.
\item Пошук найближчих номіналів E12, E24 та E96.
\item Розрахунок послідовно-паралельних резисторних мереж.
\item Визначення меж опору з урахуванням допусків.
\item Розрахунок RC-фільтрів і побудова АЧХ.
\item Розшифрування маркування резисторів і конденсаторів.
\item Розрахунок з'єднань конденсаторів, перерахунок частоти, періоду та фази.
\item Збереження історії розрахунків.
\end{enumerate}

Для структури додатка я обрав MVVM. Введення даних і результати будуть на екранах Compose, а обчислення — в окремих модулях. ViewModel зв'язуватиме екран із потрібним розрахунком. Історію зберігатиму на пристрої. Так формули можна буде перевіряти окремо та використати у версії для ПК.

\begin{figure}[H]
	\centering
	\begin{tikzpicture}[
		scale=0.88,
		transform shape,
		box/.style={draw=black!60,rounded corners=3pt,line width=0.6pt,
			align=center,inner sep=9pt,font=\small,text width=3.8cm,
			minimum height=1.8cm,fill=black!3},
		main/.style={box,text width=10.5cm,minimum height=1.3cm},
		flow/.style={<->,>=Latex,line width=0.7pt,draw=black!70}
		]
		\node[main] (ui) {\textbf{Інтерфейс · Jetpack Compose}\\Каталог інструментів, параметри та результати};
		\node[main,below=0.65cm of ui] (vm) {\textbf{ViewModel}\\Стан екрана, перевірка введення та запуск операцій};
		\node[box,below=1cm of vm] (ai) {\textbf{AI-даташити}\\PDF $\to$ локальна модель\\$\to$ картка компонента};
		\node[box,left=0.35cm of ai] (calc) {\textbf{Розрахунки}\\Компоненти, сигнали,\\датчики, RF та PCB};
		\node[box,right=0.35cm of ai] (store) {\textbf{Локальні дані}\\Історія, картки,\\PDF та налаштування};
		\node[main,below=0.7cm of ai,minimum height=1.1cm] (core) {\textbf{Спільне ядро Kotlin}\\Величини, одиниці та математичні функції};
		\draw[flow] (ui) -- (vm);
		\draw[flow] (vm.south -| calc.north) -- (calc.north);
		\draw[flow] (vm) -- (ai);
		\draw[flow] (vm.south -| store.north) -- (store.north);
		\draw[flow] (calc.south) -- (calc.south |- core.north);
	\end{tikzpicture}
	\caption{Концептуальна структура Ohmlette}
\end{figure}

\section{Git-репозиторій}
Публічний репозиторій проєкту:
\begin{center}
\small\texttt{https://github.com/mykhailo-chyzhevskyi/Ohmlette}
\end{center}
У репозиторії зберігатимуться код проєкту та README.md. Вміст README:

\begin{quote}
\textbf{Ohmlette}

Офлайн-калькулятор електронного інженера для Android.

Можливості: одиниці, стандартні номінали, резисторні мережі та допуски, RC-фільтри, маркування компонентів, з'єднання конденсаторів, історія розрахунків.

Технології: Kotlin, Jetpack Compose, MVVM.
\end{quote}

\section*{Висновок}
Під час роботи я налаштував середовище розробки й обрав тему додатка Ohmlette. Визначив, які розрахунки додам першими та як організую код. Спочатку планую реалізувати перерахунок одиниць і підбір номіналів, а потім додати інші розрахунки.
\end{document}
